\chapter{Introduction}\label{ch-intro}

\section{Algol 68}

Algol 68 was designed by Working Group 2.1 of the International Federation
for Information Processing (IFIP), the group that had earlier produced
Algol 60. The language was defined in the Report on the Algorithmic
Language Algol 68 of 1968, published in 1969~\cite{report69}; a revised definition, the
Revised Report~\cite{rr}, was published in 1975. This book describes the
language of the Revised Report, which is what all current implementations
follow. Lindsey's history of the language~\cite{history} tells how it
came about.

Algol 68 is a compiled, statically typed, imperative language. Its
designers aimed for a small number of general concepts that combine
without special cases, a property they called \term{orthogonality}. The
result is a language in which, for example, a block, an \code|if| and a
\code|case| are all expressions with a value, any type can be the element
type of an array or the result type of a function, and a function is a
value like any other.

\section{A first program}

Here is a complete program:

\example{ch01/hello}

The line after the program is its output. The program consists of a
\term{block}, the text from \code|begin| to \code|end|. The block contains
one call of the procedure \code|print|, which writes its argument to the
standard output. The argument is a list of two items in parentheses: the
string \code|"Hello, world!"| and \code|newline|, which ends the line.
That is why there are two pairs of parentheses: the outer pair belongs to
the call, the inner pair encloses the list. A call with a single item
needs only one pair, as in \code|print (newline)|.

The first line, \code|access Genie_Transput|, makes \code|print| and the
other input and output procedures available. The Revised Report includes
\code|print| in the standard library, and Algol 68 Genie provides it that
way. The standard library of ga68 provides only low-level input and
output, and a68j follows ga68. The module \code|Genie_Transput| fills the
gap: it provides the procedures of the Report that this book uses. It
comes with a68j, and its source can be compiled with ga68 as well. Every
program in this book that reads or writes begins with this line; with
Algol 68 Genie, the line is left out. Modules, a GNU extension of the
language, are described in Appendix~\ref{app-modules}.

\section{Compiling and running}

With a68j, a program is compiled and run in one step:

\begin{lstlisting}[style=shell]
$ java -jar a68j-backend.jar -run hello.a68
Hello, world!
\end{lstlisting}

With ga68, the module \code|Genie_Transput| is compiled once, and the
program is compiled and linked with it:

\begin{lstlisting}[style=shell]
$ ga68 -c genietransput.a68
$ ga68 -I. hello.a68 genietransput.o -o hello
$ ./hello
Hello, world!
\end{lstlisting}

The source of the module, \texttt{genietransput.a68}, comes with a68j.

\section{Reserved words, type names and operators}\label{sec-stropping}

Some words in an Algol 68 program are \term{reserved words}: \code|begin|,
\code|end|, \code|int|, \code|if| and some sixty others. In this book
they are printed in bold, a tradition that goes back to the Report, where
these words are written in a typeface of their own. In a source file they
are written in lower case.

How the reserved words are told apart from other words in a source file is
called \term{stropping} (\rr{stropping}). Several conventions have been
used. This book uses the default convention of ga68, which is also
understood by a68j:
\begin{itemize}
\item reserved words are written in lower case: \code|begin|, \code|int|;
\item a type name that the program declares begins with a capital letter:
  \code|Node|, \code|Point|;
\item operators that are words are written in capitals: \code|ABS|,
  \code|MOD|, or a user-defined \code|MAX|;
\item everything else, such as the names of constants, variables and
  procedures, is written in lower case, with digits and underscores:
  \code|max_value|, \code|x2|.
\end{itemize}
Algol 68 Genie uses a different convention by default, in which reserved
words, type names and word operators are written in capitals
(\code|BEGIN INT n = 1; ... END|) and names may contain spaces.

\section{A second program}

The next program computes the area and the circumference of a circle:

\example{ch01/circle}

The text between braces is a \term{comment}. The program declares three
names for values: \code|r| for the radius, and \code|area| and
\code|circumference| for the results; \code|pi| is a constant of the
standard library. The operator \code|**| raises to a power. The
procedure \code|fixed| converts a number into a string with a given number
of digits after the decimal point; the 0 asks for a string just as wide as
needed. All of this is explained in the next chapter.
